% Options for packages loaded elsewhere
\PassOptionsToPackage{unicode}{hyperref}
\PassOptionsToPackage{hyphens}{url}
\documentclass[
]{article}
\usepackage{xcolor}
\usepackage[margin=1in]{geometry}
\usepackage{amsmath,amssymb}
\setcounter{secnumdepth}{-\maxdimen} % remove section numbering
\usepackage{iftex}
\ifPDFTeX
  \usepackage[T1]{fontenc}
  \usepackage[utf8]{inputenc}
  \usepackage{textcomp} % provide euro and other symbols
\else % if luatex or xetex
  \usepackage{unicode-math} % this also loads fontspec
  \defaultfontfeatures{Scale=MatchLowercase}
  \defaultfontfeatures[\rmfamily]{Ligatures=TeX,Scale=1}
\fi
\usepackage{lmodern}
\ifPDFTeX\else
  % xetex/luatex font selection
\fi
% Use upquote if available, for straight quotes in verbatim environments
\IfFileExists{upquote.sty}{\usepackage{upquote}}{}
\IfFileExists{microtype.sty}{% use microtype if available
  \usepackage[]{microtype}
  \UseMicrotypeSet[protrusion]{basicmath} % disable protrusion for tt fonts
}{}
\makeatletter
\@ifundefined{KOMAClassName}{% if non-KOMA class
  \IfFileExists{parskip.sty}{%
    \usepackage{parskip}
  }{% else
    \setlength{\parindent}{0pt}
    \setlength{\parskip}{6pt plus 2pt minus 1pt}}
}{% if KOMA class
  \KOMAoptions{parskip=half}}
\makeatother
\usepackage{longtable,booktabs,array}
\usepackage{calc} % for calculating minipage widths
% Correct order of tables after \paragraph or \subparagraph
\usepackage{etoolbox}
\makeatletter
\patchcmd\longtable{\par}{\if@noskipsec\mbox{}\fi\par}{}{}
\makeatother
% Allow footnotes in longtable head/foot
\IfFileExists{footnotehyper.sty}{\usepackage{footnotehyper}}{\usepackage{footnote}}
\makesavenoteenv{longtable}
\usepackage{graphicx}
\makeatletter
\newsavebox\pandoc@box
\newcommand*\pandocbounded[1]{% scales image to fit in text height/width
  \sbox\pandoc@box{#1}%
  \Gscale@div\@tempa{\textheight}{\dimexpr\ht\pandoc@box+\dp\pandoc@box\relax}%
  \Gscale@div\@tempb{\linewidth}{\wd\pandoc@box}%
  \ifdim\@tempb\p@<\@tempa\p@\let\@tempa\@tempb\fi% select the smaller of both
  \ifdim\@tempa\p@<\p@\scalebox{\@tempa}{\usebox\pandoc@box}%
  \else\usebox{\pandoc@box}%
  \fi%
}
% Set default figure placement to htbp
\def\fps@figure{htbp}
\makeatother
\setlength{\emergencystretch}{3em} % prevent overfull lines
\providecommand{\tightlist}{%
  \setlength{\itemsep}{0pt}\setlength{\parskip}{0pt}}
\usepackage{bookmark}
\IfFileExists{xurl.sty}{\usepackage{xurl}}{} % add URL line breaks if available
\urlstyle{same}
\hypersetup{
  hidelinks,
  pdfcreator={LaTeX via pandoc}}

\author{}
\date{\vspace{-2.5em}}

\begin{document}

\section{🧠 Les Biais Cognitifs - Le
Jeu}\label{les-biais-cognitifs---le-jeu}

\subsection{Règles du jeu pour les
participants}\label{ruxe8gles-du-jeu-pour-les-participants}

\begin{center}\rule{0.5\linewidth}{0.5pt}\end{center}

\subsection{🎯 Objectif de la journée}\label{objectif-de-la-journuxe9e}

!

À la fin de cette journée, vous serez capable de :

\begin{itemize}
\tightlist
\item
  \textbf{Reconnaître} les principaux biais cognitifs dans votre
  quotidien professionnel
\item
  \textbf{Comprendre} comment ces biais influencent vos décisions
\item
  \textbf{Appliquer} des stratégies concrètes pour limiter leur impact
\end{itemize}

\begin{center}\rule{0.5\linewidth}{0.5pt}\end{center}

\subsection{🃏 Les 52 cartes : Qu'est-ce qu'un biais cognitif
?}\label{les-52-cartes-quest-ce-quun-biais-cognitif}

Un \textbf{biais cognitif} est un raccourci mental que notre cerveau
utilise pour prendre des décisions rapidement.

\subsubsection{Pourquoi notre cerveau fait ça
?}\label{pourquoi-notre-cerveau-fait-uxe7a}

\begin{itemize}
\tightlist
\item
  Pour \textbf{économiser de l'énergie} (penser demande beaucoup
  d'énergie !)
\item
  Pour \textbf{décider rapidement} face à trop d'informations
\item
  C'est un \textbf{mécanisme de survie} développé depuis des millénaires
\end{itemize}

\subsubsection{Le problème ?}\label{le-probluxe8me}

Ces raccourcis peuvent nous faire :

\begin{itemize}
\tightlist
\item
  \textbf{Tirer des conclusions erronées}
\item
  \textbf{Prendre de mauvaises décisions}
\item
  \textbf{Mal interpréter les situations}
\end{itemize}

\begin{quote}
{[}!IMPORTANT{]} \textbf{Point clé} : TOUT LE MONDE est sujet aux biais
cognitifs, même les experts !\\
Ce n'est pas une faiblesse, c'est le fonctionnement normal du cerveau
humain.
\end{quote}

\begin{center}\rule{0.5\linewidth}{0.5pt}\end{center}

\subsection{🎨 Les 5 catégories de
biais}\label{les-5-catuxe9gories-de-biais}

Les 52 cartes sont classées en 5 familles, chacune avec une
\textbf{couleur} et un \textbf{symbole} :

\subsubsection{1. 🧠 Mémoire \& Souvenirs
(Violet)}\label{muxe9moire-souvenirs-violet}

\textbf{Ce que c'est :} Biais qui influencent ce dont on se souvient et
comment on s'en souvient.

\textbf{Exemples :}

\begin{itemize}
\tightlist
\item
  \textbf{Effet de supériorité de l'image} : On se souvient mieux des
  images que des textes
\item
  \textbf{Effet Google} : On oublie ce qu'on peut chercher en ligne
\item
  \textbf{Effet de primauté} : On se souvient mieux du début et de la
  fin
\end{itemize}

\textbf{Impact professionnel :} Présentations, formations, communication

\begin{center}\rule{0.5\linewidth}{0.5pt}\end{center}

\subsubsection{2. 🤔 Raisonnement \& Résolution de problèmes
(Bleu)}\label{raisonnement-ruxe9solution-de-probluxe8mes-bleu}

\textbf{Ce que c'est :} Biais qui changent notre façon de raisonner et
de résoudre des problèmes.

\textbf{Exemples :}

\begin{itemize}
\tightlist
\item
  \textbf{Biais de confirmation} : On cherche ce qui confirme nos idées
\item
  \textbf{Loi de l'instrument} : ``Quand on a un marteau, tout ressemble
  à un clou''
\item
  \textbf{Biais pro-innovation} : ``Nouveau = forcément mieux''
\end{itemize}

\textbf{Impact professionnel :} Gestion de projet, innovation, analyse

\begin{center}\rule{0.5\linewidth}{0.5pt}\end{center}

\subsubsection{3. ⚡ Prise de décision \& Comportements
(Vert)}\label{prise-de-duxe9cision-comportements-vert}

\textbf{Ce que c'est :} Biais qui orientent nos choix et nos
comportements.

\textbf{Exemples :}

\begin{itemize}
\tightlist
\item
  \textbf{Biais d'ancrage} : La première info influence tout le reste
\item
  \textbf{Effet par défaut} : On choisit l'option par défaut
\item
  \textbf{Effet IKEA} : On surestime la valeur de ce qu'on a créé
\end{itemize}

\textbf{Impact professionnel :} Achats, stratégie, négociation

\begin{center}\rule{0.5\linewidth}{0.5pt}\end{center}

\subsubsection{4. 👥 Travail d'équipe \& Réunions
(Orange)}\label{travail-duxe9quipe-ruxe9unions-orange}

\textbf{Ce que c'est :} Biais qui influencent la dynamique de groupe.

\textbf{Exemples :}

\begin{itemize}
\tightlist
\item
  \textbf{Effet Dunning-Kruger} : Les incompétents surestiment leurs
  capacités
\item
  \textbf{Planification fallacieuse} : On sous-estime toujours les
  délais
\item
  \textbf{Biais d'autocomplaisance} : Succès = moi, échecs = les autres
\end{itemize}

\textbf{Impact professionnel :} Management, collaboration, réunions

\begin{center}\rule{0.5\linewidth}{0.5pt}\end{center}

\subsubsection{5. 🔍 Interviews \& Tests utilisateur
(Rose)}\label{interviews-tests-utilisateur-rose}

\textbf{Ce que c'est :} Biais qui affectent la recherche et les
entretiens.

\textbf{Exemples :}

\begin{itemize}
\tightlist
\item
  \textbf{Effet de courtoisie} : On dit ce qui est poli, pas la vérité
\item
  \textbf{Biais rétrospectif} : ``Je le savais depuis le début'' (non !)
\item
  \textbf{Biais d'observation} : Le chercheur influence l'expérience
\end{itemize}

\textbf{Impact professionnel :} Recrutement, tests, études

\begin{center}\rule{0.5\linewidth}{0.5pt}\end{center}

\subsection{🎮 Les règles
générales}\label{les-ruxe8gles-guxe9nuxe9rales}

\subsubsection{Formation d'équipe}\label{formation-duxe9quipe}

\begin{itemize}
\tightlist
\item
  La journée se joue en \textbf{équipes de 3-4 personnes}
\item
  Chaque équipe reçoit un \textbf{nom d'équipe} et une \textbf{couleur}
\item
  Les points sont comptabilisés par équipe
\end{itemize}

\subsubsection{Principe de scoring}\label{principe-de-scoring}

\begin{itemize}
\tightlist
\item
  Chaque activité rapporte des \textbf{points}
\item
  Le \textbf{tableau de score} est mis à jour en temps réel
\item
  L'équipe gagnante sera annoncée en fin de journée 🏆
\end{itemize}

\subsubsection{Esprit du jeu}\label{esprit-du-jeu}

\begin{itemize}
\tightlist
\item
  \textbf{Collaboration} : Travaillez ensemble, partagez vos idées
\item
  \textbf{Bienveillance} : Pas de jugement, on apprend tous ensemble
\item
  \textbf{Curiosité} : Posez des questions, explorez, expérimentez
\item
  \textbf{Honnêteté} : Partagez vos vraies expériences, même les erreurs
  !
\end{itemize}

\begin{center}\rule{0.5\linewidth}{0.5pt}\end{center}

\subsection{📋 Programme de la
journée}\label{programme-de-la-journuxe9e}

Présentation, formation des équipes, première expérience des biais

\subsubsection{\texorpdfstring{\textbf{Phase 1 : Découverte des
Catégories}}{Phase 1 : Découverte des Catégories}}\label{phase-1-duxe9couverte-des-catuxe9gories}

Découvrez les 52 biais et les 5 familles grâce à un quiz interactif

\subsubsection{\texorpdfstring{\textbf{Phase 2 : Chasseurs de
biais}}{Phase 2 : Chasseurs de biais}}\label{phase-2-chasseurs-de-biais}

Identifiez les biais dans des scénarios professionnels réels

\subsubsection{\texorpdfstring{\textbf{Phase 3 : Projet
Diabolique}}{Phase 3 : Projet Diabolique}}\label{phase-3-projet-diabolique}

Créez l'expérience la plus manipulatrice possible (pour comprendre les
mécanismes !)

\subsubsection{\texorpdfstring{\textbf{Phase 4 : Détox
cognitive}}{Phase 4 : Détox cognitive}}\label{phase-4-duxe9tox-cognitive}

Apprenez à identifier vos biais et découvrez des stratégies anti-biais

\subsubsection{\texorpdfstring{\textbf{Grand Quiz Final \& Proclamation
du vainqueur
🏆}}{Grand Quiz Final \& Proclamation du vainqueur 🏆}}\label{grand-quiz-final-proclamation-du-vainqueur}

\begin{center}\rule{0.5\linewidth}{0.5pt}\end{center}

\subsection{🎲 Détail des activités}\label{duxe9tail-des-activituxe9s}

\subsubsection{🔍 Activité 1 :
Découverte}\label{activituxe9-1-duxe9couverte}

\paragraph{\texorpdfstring{\textbf{Étape 1 : Classification des cartes
(15
min)}}{Étape 1 : Classification des cartes (15 min)}}\label{uxe9tape-1-classification-des-cartes-15-min}

\textbf{Ce que vous allez faire :}

\begin{itemize}
\tightlist
\item
  Votre équipe reçoit 10-12 cartes de biais cognitifs
\item
  Lisez-les ensemble et discutez-en
\item
  Classez-les par catégorie (les 5 couleurs)
\end{itemize}

\textbf{Consignes :}

\begin{enumerate}
\def\labelenumi{\arabic{enumi}.}
\tightlist
\item
  Lisez chaque carte à voix haute dans l'équipe
\item
  Essayez de comprendre chacune avec vos propres mots
\item
  Identifiez le symbole et la couleur
\item
  Regroupez les cartes par famille
\end{enumerate}

\textbf{Points} : Participation (pas de scoring compétitif sur cette
phase)

\begin{center}\rule{0.5\linewidth}{0.5pt}\end{center}

\paragraph{\texorpdfstring{\textbf{Étape 2 : Mise en
commun}}{Étape 2 : Mise en commun}}\label{uxe9tape-2-mise-en-commun}

\textbf{Ce que vous allez faire :}

\begin{itemize}
\tightlist
\item
  Chaque équipe présente 2-3 biais découverts
\item
  Expliquez avec vos propres mots ce que vous avez compris
\item
  Donnez un exemple personnel si possible
\end{itemize}

\textbf{Consignes :}

\begin{itemize}
\tightlist
\item
  Désignez un porte-parole par équipe
\item
  Soyez concis : 3-4 minutes maximum
\item
  Essayez de donner un exemple concret
\end{itemize}

\textbf{Points} : Pas de points, c'est une phase de découverte
collaborative

\begin{center}\rule{0.5\linewidth}{0.5pt}\end{center}

\paragraph{\texorpdfstring{\textbf{Étape 3 : Quiz - ``Quelle catégorie
?''}}{Étape 3 : Quiz - ``Quelle catégorie ?''}}\label{uxe9tape-3-quiz---quelle-catuxe9gorie}

\textbf{Ce que vous allez faire :}

\begin{itemize}
\tightlist
\item
  L'animateur lit des situations
\item
  Votre équipe identifie la catégorie de biais concernée
\item
  Levez votre pancarte avec la couleur correspondante
\end{itemize}

\textbf{Consignes :}

\begin{itemize}
\tightlist
\item
  Discutez rapidement en équipe (30 secondes max)
\item
  Levez votre réponse en même temps que les autres équipes
\item
  L'animateur valide ou corrige
\end{itemize}

\textbf{Barème} : 2 points par bonne réponse

\begin{center}\rule{0.5\linewidth}{0.5pt}\end{center}

\subsubsection{🎯 Activité 2 : Chasseurs de
biais}\label{activituxe9-2-chasseurs-de-biais}

\paragraph{\texorpdfstring{\textbf{Objectif}}{Objectif}}\label{objectif}

Identifier les biais cognitifs cachés dans des scénarios professionnels
réels

\paragraph{\texorpdfstring{\textbf{Comment ça marche
?}}{Comment ça marche ?}}\label{comment-uxe7a-marche}

\textbf{1. Distribution des scénarios}

\begin{itemize}
\tightlist
\item
  Votre équipe reçoit \textbf{3 scénarios professionnels}
\item
  Ce sont des situations du quotidien en entreprise
\item
  Chaque scénario contient plusieurs biais cachés
\end{itemize}

\textbf{2. Phase d'analyse} Votre mission pour chaque scénario :

\begin{itemize}
\tightlist
\item
  \textbf{Lire} attentivement le scénario
\item
  \textbf{Identifier} TOUS les biais cognitifs présents
\item
  \textbf{Utiliser vos cartes} comme aide-mémoire
\item
  \textbf{Justifier} pourquoi c'est ce biais
\item
  \textbf{Trouver un exemple} similaire dans votre expérience (bonus)
\end{itemize}

\textbf{3. Présentations}

\begin{itemize}
\tightlist
\item
  Chaque équipe présente 1 de ses scénarios
\item
  Expliquez les biais identifiés
\item
  Les autres équipes peuvent compléter ou challenger
\end{itemize}

\textbf{Conseils pour réussir :}

\begin{itemize}
\tightlist
\item
  Ne vous arrêtez pas au premier biais trouvé, cherchez plus loin !
\item
  Utilisez vos cartes : relisez-les plusieurs fois
\item
  Discutez en équipe : chacun voit des choses différentes
\item
  Pensez aux 5 catégories : y a-t-il un biais de chaque famille ?
\end{itemize}

\textbf{Barème :}

\begin{itemize}
\tightlist
\item
  \textbf{3 points} par biais correctement identifié et justifié
\item
  \textbf{2 points bonus} si vous trouvez TOUS les biais du scénario
\item
  \textbf{1 point bonus} si vous donnez un exemple concret de votre
  entreprise
\item
  \textbf{1 point bonus} si une autre équipe complète votre analyse
\end{itemize}

\begin{quote}
{[}!TIP{]} \textbf{Astuce} : Certains scénarios peuvent contenir 5 à 7
biais différents. Ne vous contentez pas des plus évidents !
\end{quote}

\begin{center}\rule{0.5\linewidth}{0.5pt}\end{center}

\subsubsection{😈 Activité 3 : Le Projet
Diabolique}\label{activituxe9-3-le-projet-diabolique}

\paragraph{\texorpdfstring{\textbf{Objectif}}{Objectif}}\label{objectif-1}

Créer le produit, service ou expérience LA PLUS MANIPULATRICE possible
en utilisant les biais cognitifs

\begin{quote}
{[}!WARNING{]} \textbf{Attention} : Cet exercice est purement
pédagogique ! L'objectif est de COMPRENDRE les mécanismes de
manipulation pour mieux s'en protéger, PAS de les appliquer réellement !
\end{quote}

\paragraph{\texorpdfstring{\textbf{Comment ça marche
?}}{Comment ça marche ?}}\label{comment-uxe7a-marche-1}

\textbf{1. Introduction et exemples}

\begin{itemize}
\tightlist
\item
  L'animateur présente des exemples de ``Dark UX'' (design manipulateur)
\item
  Vous comprenez comment les biais peuvent être utilisés pour manipuler
\end{itemize}

\textbf{2. Phase de création}

\textbf{Votre mission :} Créez un produit, service, publicité, ou
expérience qui utilise \textbf{AU MINIMUM 7 biais cognitifs} différents
pour manipuler les utilisateurs.

\textbf{Contraintes :}

\begin{itemize}
\tightlist
\item
  Utiliser au moins \textbf{1 biais de chaque catégorie} (les 5
  familles)
\item
  \textbf{Documenter} précisément quel biais est utilisé et COMMENT
\item
  Créer un support visuel (dessin, storyboard, mockup, pitch)
\end{itemize}

\textbf{Exemples de créations possibles :}

\begin{itemize}
\tightlist
\item
  Une application mobile addictive
\item
  Une publicité ultra-persuasive
\item
  Un site e-commerce manipulateur
\item
  Une expérience en magasin
\item
  Une campagne marketing
\item
  Un processus de vente
\end{itemize}

\textbf{Supports à votre disposition :}

\begin{itemize}
\tightlist
\item
  ✏️ Feuilles A3 et marqueurs
\item
  📝 Post-its
\item
  🃏 Vos 52 cartes de biais
\item
  📋 Liste complète des biais
\end{itemize}

\textbf{3. Présentations ``diaboliques''}

Chaque équipe présente :

\begin{itemize}
\tightlist
\item
  \textbf{Description} de votre création
\item
  \textbf{Démonstration de chaque biais} utilisé

  \begin{itemize}
  \tightlist
  \item
    Quel biais ?
  \item
    Comment il est appliqué ?
  \item
    Quel effet attendu ?
  \end{itemize}
\item
  \textbf{Impact} sur l'utilisateur
\end{itemize}

\textbf{Barème :}

\begin{itemize}
\tightlist
\item
  \textbf{2 points} par biais correctement appliqué et justifié
\item
  \textbf{3 points bonus} pour l'équipe la plus ``diabolique'' (vote des
  participants)
\item
  \textbf{3 points bonus} pour l'équipe la plus créative (vote de
  l'animateur)
\end{itemize}

\textbf{Exemples pour vous inspirer :}

\textbf{Exemple 1 : Site de réservation d'hôtel}

\begin{itemize}
\tightlist
\item
  \textbf{Biais d'ancrage} : ``Prix initial 300€'' barré, puis ``150€''
\item
  \textbf{Effet par défaut} : Option ``assurance annulation'' pré-cochée
\item
  \textbf{Effet de groupe} : ``234 personnes regardent cet hôtel EN CE
  MOMENT''
\item
  \textbf{Actualisation hyperbolique} : ``Réservez maintenant : -20\%
  expire dans 15 min''
\item
  \textbf{Effet de rareté} : ``Plus que 1 chambre disponible !''
\item
  \textbf{Biais d'autorité} : ``Recommandé par Le Figaro''
\item
  \textbf{Effet de supériorité de l'image} : Photos magnifiques (vs
  réalité moyenne)
\end{itemize}

\textbf{Exemple 2 : Application de fitness}

\begin{itemize}
\tightlist
\item
  \textbf{Effet IKEA} : Personnalisation du programme (``VOTRE plan'')
\item
  \textbf{Biais d'unité} : ``Vous avez commencé la séance, finissez-la
  !''
\item
  \textbf{Effet de groupe} : ``12 000 membres ont fait cet exercice
  aujourd'hui''
\item
  \textbf{Planification fallacieuse} : ``Objectif : -10kg en 1 mois !''
  (irréaliste)
\item
  \textbf{Biais de confirmation} : Ne montre que les success stories
\item
  \textbf{Actualisation hyperbolique} : Récompenses immédiates (badges)
\item
  \textbf{Effet de simple exposition} : Notifications quotidiennes
\end{itemize}

\begin{quote}
{[}!CAUTION{]} \textbf{Rappel éthique} : Une fois que vous savez comment
fonctionne la manipulation :

\begin{itemize}
\tightlist
\item
  Protégez-vous en tant que consommateur
\item
  Protégez vos utilisateurs en tant que professionnel
\item
  Utilisez ces connaissances de manière éthique !
\end{itemize}
\end{quote}

\begin{center}\rule{0.5\linewidth}{0.5pt}\end{center}

\subsubsection{🛡️ Activité 4 : Détox
Cognitive}\label{activituxe9-4-duxe9tox-cognitive}

\paragraph{\texorpdfstring{\textbf{Objectif}}{Objectif}}\label{objectif-2}

Apprendre à identifier VOS propres biais et découvrir des stratégies
pour les limiter

\paragraph{\texorpdfstring{\textbf{Partie 1 : Mon profil cognitif
}}{Partie 1 : Mon profil cognitif }}\label{partie-1-mon-profil-cognitif}

\textbf{Ce que vous allez faire :}

\begin{itemize}
\tightlist
\item
  Répondre \textbf{individuellement} à un questionnaire de 20 situations
\item
  Identifier vos biais dominants
\item
  Comprendre vos patterns de pensée
\end{itemize}

\textbf{Consignes :}

\begin{itemize}
\tightlist
\item
  Soyez \textbf{honnête} avec vous-même (ce n'est pas noté !)
\item
  Choisissez la réponse qui correspond à ce que vous faites
  \textbf{VRAIMENT}, pas ce que vous ``devriez'' faire
\item
  Il n'y a pas de bonne ou mauvaise réponse
\end{itemize}

\textbf{Résultat :} Vous identifiez vos 3-5 biais les plus fréquents.
Cette prise de conscience est le premier pas pour les limiter !

\begin{center}\rule{0.5\linewidth}{0.5pt}\end{center}

\paragraph{\texorpdfstring{\textbf{Partie 2 : La boîte à outils
anti-biais}}{Partie 2 : La boîte à outils anti-biais}}\label{partie-2-la-bouxeete-uxe0-outils-anti-biais}

\textbf{Ce que vous allez apprendre :} L'animateur vous présente 7
stratégies concrètes pour limiter l'impact des biais dans vos décisions
professionnelles.

\textbf{Les 7 stratégies :}

\textbf{1. ⏱️ Ralentir la décision}

\begin{itemize}
\tightlist
\item
  Ne décidez pas sous pression
\item
  Technique : Attendre 24h pour les décisions importantes
\end{itemize}

\textbf{2. 🔄 Chercher la contradiction}

\begin{itemize}
\tightlist
\item
  Activement chercher ce qui contredit votre avis
\item
  Technique : ``Et si le contraire était vrai ?''
\end{itemize}

\textbf{3. 📚 Diversifier les sources}

\begin{itemize}
\tightlist
\item
  Ne pas se fier à une seule source
\item
  Technique : Règle des 3 sources minimum
\end{itemize}

\textbf{4. ✅ Utiliser des checklists}

\begin{itemize}
\tightlist
\item
  Structurer les décisions importantes
\item
  Technique : Grille de décision avec critères objectifs
\end{itemize}

\textbf{5. 👥 Décisions collectives structurées}

\begin{itemize}
\tightlist
\item
  Utiliser des méthodes qui limitent les biais de groupe
\item
  Technique : Brainstorming silencieux, vote anonyme
\end{itemize}

\textbf{6. 🧘 Attention au contexte}

\begin{itemize}
\tightlist
\item
  Ne pas décider si faim, fatigue, stress, émotion forte
\item
  Technique : Auto-questionnement ``Suis-je en état de décider ?''
\end{itemize}

\textbf{7. 📖 Formation continue}

\begin{itemize}
\tightlist
\item
  Rester vigilant et partager avec l'équipe
\item
  Technique : ``Biais du mois'' en réunion d'équipe
\end{itemize}

\begin{center}\rule{0.5\linewidth}{0.5pt}\end{center}

\paragraph{\texorpdfstring{\textbf{Partie 3 : Plan d'action en
équipe}}{Partie 3 : Plan d'action en équipe}}\label{partie-3-plan-daction-en-uxe9quipe}

\textbf{Ce que vous allez faire :} Créer un plan concret pour appliquer
ce que vous avez appris dans VOTRE contexte professionnel

\textbf{Consignes :}

\begin{enumerate}
\def\labelenumi{\arabic{enumi}.}
\tightlist
\item
  Choisissez \textbf{1 situation professionnelle récurrente} dans votre
  entreprise

  \begin{itemize}
  \tightlist
  \item
    Exemples : réunions projet, recrutement, gestion conflit, choix
    fournisseur
  \end{itemize}
\item
  Identifiez les \textbf{biais qui peuvent affecter} cette situation
\item
  Proposez \textbf{3 stratégies concrètes} pour limiter ces biais
\item
  Définissez \textbf{comment mesurer} si ça fonctionne
\end{enumerate}

\textbf{Livrable :} Une fiche action que vous pourrez ramener dans votre
entreprise

\textbf{Format de la fiche :}

\begin{verbatim}
SITUATION : [Ex: Réunion de lancement de projet]

BIAIS IDENTIFIÉS :
- Biais 1 : [Ex: Effet de groupe - tout le monde suit le chef]
- Biais 2 : [Ex: Loi de l'instrument - on refait toujours pareil]
- Biais 3 : [Ex: Planification fallacieuse - on sous-estime les délais]

STRATÉGIES CONCRÈTES :
1. [Ex: Brainstorming silencieux avant discussion]
2. [Ex: Challenger systématiquement "pourquoi cette méthode ?"]
3. [Ex: Multiplier les estimations initiales par 1,5]

INDICATEURS DE SUCCÈS :
- [Ex: Au moins 2 idées nouvelles par réunion]
- [Ex: Respect des délais dans 80% des cas]
\end{verbatim}

\begin{center}\rule{0.5\linewidth}{0.5pt}\end{center}

\subsubsection{🏆 Activité 5 : Grand Quiz
Final}\label{activituxe9-5-grand-quiz-final}

\paragraph{\texorpdfstring{\textbf{Format}}{Format}}\label{format}

Quiz de 20 questions style ``Qui veut gagner des millions''

\paragraph{\texorpdfstring{\textbf{Types de
questions}}{Types de questions}}\label{types-de-questions}

\textbf{QCM} : Connaissances sur les biais

\begin{itemize}
\tightlist
\item
  \emph{Exemple : ``L'effet IKEA signifie : A) On aime les meubles B) On
  surestime ce qu'on a créé C) On sous-estime les délais''}
\end{itemize}

\textbf{Vrai / Faux} : Idées reçues

\begin{itemize}
\tightlist
\item
  \emph{Exemple : ``Les experts sont immunisés contre les biais
  cognitifs'' → FAUX}
\end{itemize}

\textbf{Catégorisation} : Associer biais et famille

\begin{itemize}
\tightlist
\item
  \emph{Exemple : ``Le biais de confirmation appartient à quelle
  catégorie ?''}
\end{itemize}

\textbf{Situations} : Identifier le biais dans un mini-cas

\begin{itemize}
\tightlist
\item
  \emph{Exemple : ``Marc fait confiance à Google plutôt qu'à sa mémoire.
  Quel biais ?'' → Effet Google}
\end{itemize}

\textbf{Application} : Solutions anti-biais

\begin{itemize}
\tightlist
\item
  \emph{Exemple : ``Comment contrer le biais d'ancrage ? A) Ignorer la
  première info B) Chercher plusieurs références C) Décider vite''}
\end{itemize}

\paragraph{\texorpdfstring{\textbf{Barème}}{Barème}}\label{baruxe8me}

\begin{itemize}
\tightlist
\item
  1 point par bonne réponse
\item
  5 points bonus pour l'équipe gagnante du quiz
\item
  Le score s'ajoute au score global de la journée
\end{itemize}

\begin{center}\rule{0.5\linewidth}{0.5pt}\end{center}

\subsection{🏅 Scoring final et
vainqueur}\label{scoring-final-et-vainqueur}

\subsubsection{\texorpdfstring{\textbf{Récapitulatif des points de la
journée}}{Récapitulatif des points de la journée}}\label{ruxe9capitulatif-des-points-de-la-journuxe9e}

\begin{longtable}[]{@{}ll@{}}
\toprule\noalign{}
Activité & Points maximum \\
\midrule\noalign{}
\endhead
\bottomrule\noalign{}
\endlastfoot
Quiz ``Quelle catégorie ?'' & \textasciitilde10 points \\
Chasseurs de biais & \textasciitilde50 points \\
Projet Diabolique & \textasciitilde20 points \\
Quiz final & \textasciitilde25 points \\
\textbf{TOTAL} & \textbf{\textasciitilde105 points} \\
\end{longtable}

\subsubsection{\texorpdfstring{\textbf{L'équipe gagnante remporte
:}}{L'équipe gagnante remporte :}}\label{luxe9quipe-gagnante-remporte}

\begin{itemize}
\tightlist
\item
  🏆 Le titre de \textbf{``Chasseurs de biais 2026''}
\item
  🎖️ Un certificat d'excellence
\item
  🧠 La fierté d'avoir dominé les biais cognitifs !
\end{itemize}

\begin{center}\rule{0.5\linewidth}{0.5pt}\end{center}

\subsection{💡 Conseils pour profiter au maximum de la
formation}\label{conseils-pour-profiter-au-maximum-de-la-formation}

\subsubsection{\texorpdfstring{\textbf{1. Participez
activement}}{1. Participez activement}}\label{participez-activement}

\begin{itemize}
\tightlist
\item
  Posez des questions
\item
  Partagez vos expériences
\item
  N'ayez pas peur de vous tromper : on apprend par l'erreur !
\end{itemize}

\subsubsection{\texorpdfstring{\textbf{2. Soyez
curieux}}{2. Soyez curieux}}\label{soyez-curieux}

\begin{itemize}
\tightlist
\item
  Explorez toutes les cartes, pas seulement celles de vos scénarios
\item
  Cherchez à comprendre le ``pourquoi'' derrière chaque biais
\item
  Faites des liens avec votre quotidien
\end{itemize}

\subsubsection{\texorpdfstring{\textbf{3. Collaborez avec votre
équipe}}{3. Collaborez avec votre équipe}}\label{collaborez-avec-votre-uxe9quipe}

\begin{itemize}
\tightlist
\item
  Écoutez les points de vue différents
\item
  Chacun voit des choses que les autres ne voient pas
\item
  La diversité d'opinion est une force contre les biais !
\end{itemize}

\subsubsection{\texorpdfstring{\textbf{4. Prenez des
notes}}{4. Prenez des notes}}\label{prenez-des-notes}

\begin{itemize}
\tightlist
\item
  Notez les biais qui vous parlent le plus
\item
  Écrivez des exemples personnels pour mieux mémoriser
\item
  Gardez votre fiche action pour l'appliquer ensuite
\end{itemize}

\subsubsection{\texorpdfstring{\textbf{5. Amusez-vous
!}}{5. Amusez-vous !}}\label{amusez-vous}

\begin{itemize}
\tightlist
\item
  C'est un jeu, avant tout
\item
  L'apprentissage est plus efficace quand c'est ludique
\item
  Profitez des moments de groupe et d'échanges
\end{itemize}

\begin{center}\rule{0.5\linewidth}{0.5pt}\end{center}

\subsection{📝 Ce que vous repartirez
avec}\label{ce-que-vous-repartirez-avec}

\subsubsection{\texorpdfstring{\textbf{En fin de journée, vous aurez
:}}{En fin de journée, vous aurez :}}\label{en-fin-de-journuxe9e-vous-aurez}

✅ \textbf{Connaissance} des 52 biais cognitifs et des 5 catégories\\
✅ \textbf{Exemples concrets} de biais dans des situations
professionnelles\\
✅ \textbf{Profil cognitif} : connaissance de vos biais dominants\\
✅ \textbf{Boîte à outils} : 7 stratégies anti-biais applicables
immédiatement\\
✅ \textbf{Plan d'action} personnalisé pour votre contexte
professionnel\\
✅ \textbf{Supports de formation} : fiche récapitulative, carte mentale,
bibliographie

\begin{center}\rule{0.5\linewidth}{0.5pt}\end{center}

\subsection{🎓 Et après la formation ?}\label{et-apruxe8s-la-formation}

\subsubsection{\texorpdfstring{\textbf{Pour continuer l'apprentissage
:}}{Pour continuer l'apprentissage :}}\label{pour-continuer-lapprentissage}

\textbf{📚 Lectures recommandées}

\begin{itemize}
\tightlist
\item
  ``Thinking, Fast and Slow'' - Daniel Kahneman
\item
  ``Nudge'' - Richard Thaler
\item
  ``Predictably Irrational'' - Dan Ariely
\end{itemize}

\textbf{🌐 Ressources en ligne}

\begin{itemize}
\tightlist
\item
  Sites interactifs sur les biais cognitifs
\item
  Vidéos TED Talks
\item
  Podcasts sur l'économie comportementale
\end{itemize}

\textbf{🏢 En entreprise}

\begin{itemize}
\tightlist
\item
  Partagez vos apprentissages avec votre équipe
\item
  Mettez en place un rituel ``biais du mois''
\item
  Créez une culture de décision réfléchie
\end{itemize}

\textbf{🔄 Pratique quotidienne}

\begin{itemize}
\tightlist
\item
  Identifiez 1 biais par jour dans votre vie
\item
  Testez les stratégies anti-biais
\item
  Partagez vos découvertes
\end{itemize}

\begin{center}\rule{0.5\linewidth}{0.5pt}\end{center}

\subsection{❓ FAQ}\label{faq}

\subsubsection{\texorpdfstring{\textbf{Q : Est-ce que je peux vraiment
éliminer mes biais
?}}{Q : Est-ce que je peux vraiment éliminer mes biais ?}}\label{q-est-ce-que-je-peux-vraiment-uxe9liminer-mes-biais}

\textbf{R :} Non, et ce n'est pas le but ! Les biais font partie du
fonctionnement normal du cerveau. L'objectif est de les
\textbf{reconnaître} et de les \textbf{limiter} quand ils peuvent nuire
à vos décisions importantes.

\subsubsection{\texorpdfstring{\textbf{Q : Avoir des biais, c'est être
stupide
?}}{Q : Avoir des biais, c'est être stupide ?}}\label{q-avoir-des-biais-cest-uxeatre-stupide}

\textbf{R :} Absolument pas ! Même les prix Nobel ont des biais. C'est
un mécanisme d'économie cognitive universel. La stupidité serait de nier
qu'on en a (c'est le biais de l'angle mort !).

\subsubsection{\texorpdfstring{\textbf{Q : Comment puis-je utiliser ces
connaissances au travail
?}}{Q : Comment puis-je utiliser ces connaissances au travail ?}}\label{q-comment-puis-je-utiliser-ces-connaissances-au-travail}

\textbf{R :} Dans plein de situations : recrutement, gestion de projet,
réunions, vente, marketing, innovation, résolution de conflits, etc.
Vous verrez rapidement où les appliquer !

\subsubsection{\texorpdfstring{\textbf{Q : C'est éthique d'utiliser les
biais pour influencer
?}}{Q : C'est éthique d'utiliser les biais pour influencer ?}}\label{q-cest-uxe9thique-dutiliser-les-biais-pour-influencer}

\textbf{R :} Ça dépend du contexte et de l'intention.
\textbf{Influencer} positivement (nudge éthique) ≠ \textbf{Manipuler}
(dark UX). La ligne est fine, d'où l'importance de cette formation !

\subsubsection{\texorpdfstring{\textbf{Q : Puis-je utiliser les cartes
après la formation
?}}{Q : Puis-je utiliser les cartes après la formation ?}}\label{q-puis-je-utiliser-les-cartes-apruxe8s-la-formation}

\textbf{R :} Oui ! Vous pouvez les télécharger, les imprimer, les
utiliser avec vos équipes (licence Creative Commons).

\begin{center}\rule{0.5\linewidth}{0.5pt}\end{center}

\subsection{✨ Message final}\label{message-final}

Les biais cognitifs sont fascinants : ils révèlent comment notre cerveau
fonctionne réellement, pas comment on croit qu'il fonctionne.

Cette journée est une opportunité unique de :

\begin{itemize}
\tightlist
\item
  \textbf{Mieux vous connaître}
\item
  \textbf{Prendre de meilleures décisions}
\item
  \textbf{Comprendre les autres} (qui ont les mêmes biais que vous !)
\item
  \textbf{Développer votre esprit critique}
\end{itemize}

Profitez-en, amusez-vous, et rappelez-vous : \textbf{nous sommes TOUS
biaisés} ! 🧠

\begin{center}\rule{0.5\linewidth}{0.5pt}\end{center}

\textbf{Bonne formation et que le meilleur chasseur de biais gagne ! 🎯}

\end{document}
